% =========================================================
% CSC311 Final Project Report Template
% Single column, 12pt, 8 pages max.
% =========================================================

\documentclass[12pt]{article}

% ---------- Packages ----------
\usepackage[margin=1in]{geometry}
\usepackage{parskip}         % blank line between paragraphs, no indent
\usepackage{graphicx}
\usepackage{booktabs}
\usepackage{amsmath, amssymb}
\usepackage{siunitx}
\usepackage[hidelinks]{hyperref}
\usepackage{caption}
\usepackage{enumitem}
\usepackage{xcolor}
\usepackage{XCharter}
\usepackage[xcharter]{newtxmath}
\usepackage[scaled=0.96]{inconsolata}

%--------- Headers and footers ------
\usepackage{fancyhdr}
\pagestyle{fancy}
\setlength{\headheight}{14.49998pt}
\fancyhf{}
\fancyhead[L]{CSC311 Project Report}
\fancyhead[R]{Group Members}
\fancyfoot[C]{\thepage}

% ---------- Captions & Lists ----------
\captionsetup{font=small, labelfont=bf}
\setlist[itemize]{noitemsep, topsep=2pt}
\setlist[enumerate]{noitemsep, topsep=2pt}

% ---------- Helpful macros ----------
\newcommand{\bestmodel}{\textit{[Best Model Name]}}
\newcommand{\projacc}{\textbf{[XX.XX\%]}}

% ---------- Safe figure helper ----------
\makeatletter
\newcommand{\maybeincludegraphics}[3]{%
  \begin{figure}[h]
    \centering
    \IfFileExists{#1}{%
      \includegraphics[width=0.7\linewidth]{#1}%
    }{%
      \fbox{%
        \parbox[c][0.28\textheight][c]{0.7\linewidth}{%
          \centering
          \vspace{0.5em}
          \textit{[Placeholder for \texttt{#1}]}\\
          Add your figure here.
          \vspace{0.5em}
        }%
      }%
    }%
    \caption{#2}
    \label{#3}
  \end{figure}%
}
\makeatother

\usepackage{xcolor}
\definecolor{royalblue}{rgb}{0.0, 0.14, 0.4}
\definecolor{royalblue(web)}{rgb}{0.25, 0.41, 0.88}
\usepackage{tcolorbox}
\tcbuselibrary{breakable}
\usepackage{comment}
\usepackage{tikz}
\usetikzlibrary{arrows.meta}


% =========================================================
% Title Page
% =========================================================
\begin{document}

\begin{titlepage}
    \centering
    \vspace*{3cm}

    {\Large {CSC311: Introduction to Machine Learning} \par}
    \vspace{1cm}
    
    {\Large {Project Report} \par}
    \vspace{1cm}
    
    {\large \textit{Optional Creative Project Title} \par}
    \vspace{3cm}
    
    {\large Submitted by \par}
    
    \vspace{0.5cm}
    {\large Group Number}
    
    \vspace{0.5cm}
    {\large Student 1, Student 2, Student 3, Student 4 \par}
    
    \vfill
    {\large University of Toronto \\ Fall 2026 \par}
\end{titlepage}
\clearpage
\setcounter{page}{1}

% ================================================
% 2. Executive Summary
% ================================================
\section{Executive Summary}

This section should be at most 1/4 page (1 paragraph).

\begin{itemize}
  \item Summarize the model families you explored (at least three).  
  \item State your best-performing model and its projected accuracy.  
  \item Provide a high-level justification of why this model outperformed the others.  
\end{itemize}

\section{Data Exploration}\label{sec:data-exploration}

This section of your report should be at most 2 pages. 

The goal of the Data Exploration step is to understand your dataset well enough to decide how to prepare it and which models to try.

Examine the type and distribution of every feature. Use your findings to decide how to pre-process each feature, including normalizing, encoding, and vectorizing the text.

Look for patterns in the data, such as linear or nonlinear relationships between the features and the label. These patterns help you decide which model families are suitable.

Determine which features appear predictive, so that you can justify the models you choose in the Methodology section.

Make sure to write down all of this reasoning. State which transformations you applied and why, and explain what each plot and table reveals.

Many students in previous semesters found this section difficult, because most of data exploration is not taught explicitly in class. Appendix~\ref{app:data-exploration} gives a list of guiding questions to walk you through the process. The questions point out the common issues you should look for and the standard strategies for handling them.

\section{Methodology}\label{sec:methodology}

This section of your report should be at most 3 pages.

The goal of the Methodology step is to describe how you trained, tuned, and evaluated your models in enough detail that a reader could reproduce your results.

Present the three model families you used and explain why each one suits the dataset, drawing on what you found during data exploration. Do not jump straight to a deep network, because a simpler model may perform just as well and is easier to interpret.

Make sure to describe how you trained each model so that your results can be reproduced. State the optimizer (e.g., SGD or Adam), the learning rate schedule if you used one, and any regularization or early stopping. Briefly name the libraries you used and describe any custom code you wrote.

Explain how you split the data into training and validation sets, or how you applied cross-validation, so that the reader can trust your comparisons. Tune hyperparameters on a validation set that is separate from the training data. Never use the test set for training, feature engineering, or tuning.

Summarize your hyperparameter tuning by stating which hyperparameters you varied, how you chose the values to try, and which metric you used to compare them. An exhaustive search is not required, but you should show that your choices are reasonable.

Be sure to tune every model family before you compare them. Choosing the best model using default hyperparameters and then tuning only that model is unfair to the other models, which might improve just as much with tuning. This was the most common mistake we saw when grading reports in previous semesters.

Describe the evaluation metrics you used and justify each one. Report at least one metric beyond accuracy, such as precision, recall, or F1, because accuracy alone can hide poor performance on one of the classes. Make sure every table and plot has appropriate labels and an explanation of what it shows.

Optionally, you may describe any additional techniques you tried, such as feature engineering, clustering, or PCA.

Appendix~\ref{app:methodology} gives a list of guiding questions to work through before you write this section.

% ===========================================
% 5. Results
% ===========================================
\section{Results}\label{sec:results}

This section should be at most 1 page.

The goal of the Results step is to show how well your models performed, explain which model you chose, and predict how it will do on the unseen test set.

Report the results of every model using the evaluation metrics you described in the Methodology section. Present them in tables or plots that make the comparison easy to read.

Compare performance across the three model families, and state clearly which model you chose as your final model. This must be the model implemented in your prediction script.

Analyze the errors your models make. Describe the common misclassifications or weaknesses you observed, and include a confusion matrix or a few illustrative examples if they help.

Estimate your final model's accuracy on the unseen test set as a single number. A range will earn no marks. Make sure to support your estimate with empirical evidence, such as validation results, the stability of your cross-validation scores, learning curves, or consistency across random seeds or folds.

% ================================================
% 8. Contributions and Learning
% ================================================
\section{Contributions and Learning}

This section should contain a list of 3–4 bullets, one for each student. Each student should write at most three sentences describing their main contributions and the most important lesson learned.

% ==================================================
% 9. References
% ==================================================
\section{References}

Use references whenever appropriate (e.g., when you consult textbooks, research papers, online tutorials, or external code). Cite in the text using \verb|\cite{}|. For example, let's cite the two references provided in the \texttt{.bib} file \cite{bishop2006prml,hastie2009esl}.

\bibliographystyle{plain}
\bibliography{references}

\newpage
\appendix
% Extra space between the subtitle groups and between the questions in the appendices.
\setlist[enumerate,1]{itemsep=1em}
\setlist[enumerate,2]{itemsep=0.6em}
% ==================================================
% Appendix: guiding questions for Data Exploration
% ==================================================
\section{Guiding Questions for Data Exploration}\label{app:data-exploration}

These questions support Section~\ref{sec:data-exploration}. Work through them in order, and record your answers for our dataset in that section rather than here. This appendix is part of the template instructions, so comment it out before you submit, as you do with the rest of the instructions.

%Define Prediction Task
%
%\begin{enumerate}
%
%\item Is the task classification or regression? If it is a classification problem, what is the number of classes? 
%
%\end{enumerate}

\begin{enumerate}

\item Determine Input to Model 

\begin{enumerate}

\item For each example, what features do we have? What are the types of the features (numerical, categorical, or text)? 

\item In the released data set, each row has a unique ID identifying which student provided the data. Why should we not use this unique ID as a feature?

% What is the distribution of each feature over all the examples? Is the feature highly skewed? Is there enough variation? Are there outliers? 

\end{enumerate}

\item Data Splitting

\begin{enumerate}

\item Why is it a good idea to split the provided data into a training set and a test set? 

\item Why is it a good idea to do the training and test set splitting based on the unique ID identifying each respondent? 

\item What is the role of the validation set and why do we need this third split?

\item Why should we only perform data exploration on the training set and not on the validation or the test sets?

\end{enumerate}

\item Feature Distributions

\smallskip
One reason for exploring the distributions of features is to get to know the data before building any model. This exploration can help us catch common issues such as outliers, skewed distributions, class imbalance, missing values, and large difference in feature scales.
   
\smallskip
\begin{enumerate}

\item What should we do when exploring the distributions of the features?

\item Suppose a feature takes the same value for almost every respondent. Why is this feature worth finding, and what would you do about it?

 \item Suppose two features are highly correlated, such as two Likert questions that respondents tend to answer the same way. How would you test whether two features are highly correlated, and why does a highly correlated pair cause trouble?

\end{enumerate}

\item Class Imbalance

\begin{enumerate}

\item What is class imbalance? Why is class imbalance problematic? 

\item Does the dataset have a class imbalance issue? If so, how severe?

\item What are some common strategies to address class imbalance? 

\end{enumerate}

\item Missing Values

\begin{enumerate}
 
\item First, determine if the dataset contains missing values. Report the percentage of values missing for each feature in a table. A missing value can take many forms such as a null value, an empty string, an empty list, ``n/a'', or ``-''. Normalize all of these to a single missing marker before you count, or you will underestimate how much is missing.
\bigskip

\item Next, for each feature with a substantial missing rate, try to determine why the values are likely missing. Different explanations may require different fixes. 

What are some possible explanations for missing values?

\bigskip

\item
Since you never see the missing value, it may not be possible to verify which explanation holds. However, it would be useful to determine whether the missingness of a feature is related to the label or not. If the missingness is related to the label, a missing value itself provides a predictive signal. What is a common strategy for determining this? 

\item What are some common strategies to handle missing values? One is called imputation. 

\item When is it appropriate to use each strategy above to handle missing values?

\item 
Suppose that you decide to impute the median value for a numeric feature. Why is it problematic to do the imputation on the validation and test sets?

\item Finally, decide how missing values should be handled for each feature in our dataset, and record your decision.

\end{enumerate}

\item Outliers

\begin{enumerate}

\item What observation might cause you to suspect that certain values are outliers? 

\item How do I determine whether a feature value should be considered an outlier?

\item What are some strategies to handle outliers?

\end{enumerate}

\item Feature Engineering

\begin{enumerate}

\item How can we convert the responses to a multi-select question into a feature? 

\item The survey has a few Likert-scale questions that ask respondents to indicate their level of agreement to a statement using an ordered set of options (e.g., from Strongly Disagree to Strongly Agree). We can convert the Likert-scale responses to features by modeling them as ordinal values or one-hot encoded categories. What are some advantages and disadvantages of each approach?

\item At least one free-text question asks for a number. How would you turn such a text response into a number, and when would you declare a response unparseable?

\item What are some strategies to convert the free-text responses to numeric features?

\item One common technique for processing the free-text responses is TF-IDF vectorization. Explain this technique in detail. Why is this technique suitable for processing the free-text data? 

\item How should we handle a word that appears at prediction time but never in the training split?

\item After encoding all the features, you recorded that there are 1100 features and 300 examples. How would the ratio between these two inform your model choices? Why would you choose a simpler model with strong regularization over a larger model?

\end{enumerate}

\item Feature Scaling

\begin{enumerate}

\item Two common choices for feature scaling are standardization and min-max scaling. Explain each in one sentences. Why do we need to perform feature scaling for our dataset, which has Likert-scale, multi-hot encoded, and word count features? 

\item Is each model below affected by feature scaling? Briefly explain.
\begin{itemize}
\item K-nearest neighbours
\item Decision trees and random forest
\item Softmax regression trained using gradient descent
\item L2 Regularization
\item Naive Bayes
\end{itemize}

\item You decide to standardize a feature. Which split do you use to compute the mean and standard deviation, and why?

\end{enumerate}

\item Which Features Are Predictive

\smallskip
Next, we would like to determine whether a feature is useful by looking at the feature together with the label. This analysis could also tell you which model families are worth trying and help you justify your model choices in the Methodology section.
\smallskip

\begin{enumerate}

\item How do you examine the relationship between a Likert-scale or numeric feature and the label?

\item How do you examine the relationship between a multi-select or free-text feature and the label?

\item Why should you only do this analysis on the training split and not on the validation or test splits?

\item How can you delete whether a feature is worth keeping by comparing the difference in mean feature values between classes with the spread of feature values within a class?

\smallskip
Consider a Likert-scale feature on a scale of 1 to 5. Is the feature worth keeping if two classes have means 3.5 and 3.8 while the standard deviation of answers within a class is 1.2? What about if two classes have means 2.1 and 4.3 while the standard deviation of answers within a class is 0.6?

\item What does the shape of the relationship tell you about which model families to try? Consider the three scenarios below. Discuss the model families you would try for each one.
\begin{enumerate}
\item The classes separate well along one feature at a threshold.
\item The middle values of a Likert feature belong to one class while both extremes belong to another.
\item After vectorizing the free-text responses, you have several hundred word-count features. No single word separates the classes, but each word in a response carries a small amount of evidence.
\end{enumerate}

\item Finally, write down the features you expect to be the most predictive, and name the plot or the number that led you to each one.

\end{enumerate}

\end{enumerate}

% ==================================================
% Appendix: guiding questions for Methodology
% ==================================================
\section{Guiding Questions for Methodology}\label{app:methodology}

These questions support Section~\ref{sec:methodology}. Work through them in order, and record your answers for our dataset in that section rather than here. This appendix is part of the template instructions, so comment it out before you submit, as you do with the rest of the instructions.

\begin{enumerate}

\item Choosing Model Families

\begin{enumerate}

\item Based on the size of the data set, should we consider fine-tuning a large language model or starting with simpler models? What are some advantages and disadvantages of each approach?

\item Which findings from your data exploration would support choosing softmax regression? Which findings would support choosing a decision tree instead?

\end{enumerate}

\item Training Details

\begin{enumerate}

\item Suppose you train a multilayer perceptron. Which choices would another student need to know to reproduce your results exactly?

\item Your softmax regression model reaches 95\% accuracy on the training set but only 70\% on the validation set. What does this gap tell you, and which regularization or early stopping strategies could you try?

\end{enumerate}

\item Validation

\begin{enumerate}

\item What is cross validation? If we want to use cross validation, why should we use grouped cross validation, grouped by the unique ID?

\item Pre-processing steps such as normalization and text vectorization learn values from the data, like a mean or a vocabulary. Which part of the data should these values be computed from during cross validation, and why?

\end{enumerate}

\item Hyperparameter Tuning

\begin{enumerate}

\item Which hyperparameters would you tune for a decision tree, and how would you choose the range of values to try?

\item You try maximum depths from 2 to 10 for a decision tree, and the validation accuracy is highest at a depth of 10. Is your search finished? What would you do next?

\item After choosing your final model and its hyperparameters, should you retrain it on the training and validation sets combined? What do you gain, and what do you lose?

\end{enumerate}

\item Comparing Model Families Fairly

\begin{enumerate}

\item With default hyperparameters, softmax regression reaches 78\% validation accuracy and a decision tree reaches 72\%. You tune softmax regression up to 82\% and choose it as your final model. What is wrong with this process, and how would you fix it?

\item Across five folds, model A has a mean validation accuracy 1\% higher than model B, but the accuracy of each model varies by about 3\% from fold to fold. Can you conclude that model A is better? What else could help you decide?

\end{enumerate}

\item Evaluation Metrics

\begin{enumerate}

\item Suppose one class makes up 60\% of the training set. What accuracy does a model that always predicts this class achieve? Which metric would reveal that this model is useless?

\item How do you compute precision, recall, and F1 for a 3-class problem? What is the difference between macro-averaging and weighted averaging, and when would you prefer each one?

\end{enumerate}

\item Additional Techniques

\begin{enumerate}

\item Suppose you apply PCA to several hundred word-count features. How would you decide how many components to keep, and which part of the data should you fit PCA on?

\item How could you use a clustering method such as k-means to explore the free-text responses? How would you tell whether the clusters improve your predictions?

\end{enumerate}

\end{enumerate}

\end{document}
